% !TEX root = AImath.v4.tex

%认知模式：拉回“智能的构建方式”，强调多路径融合（神经-符号-因果）

\chapter{符号、神经与因果：智能的三种建构方式 \texorpdfstring{\\}{ } 三条道路，一场融合}

\section*{回到起点：我们是如何理解“智能”的？}

在前面的章节中，我们讨论了深度学习的奇迹与幻觉，也看到了它所触及的计算边界。或许你也开始思考一个问题：如果神经网络不是“智能”的终极方式，那我们还有其他的建构路径吗？

答案是肯定的。人类在追寻人工智能的历程中，曾走过三条主要路线：

\begin{itemize}
\item \textbf{符号主义}（Symbolism）：将智能理解为符号的操控与逻辑规则的运用；
\item \textbf{连接主义}（Connectionism）：认为智能是神经元网络中的权重分布与信息激活；
\item \textbf{因果主义}（Causalism）：主张智能不仅要“预测”，更要“解释”与“推理因果”。
\end{itemize}

这三种认知模式，分别模仿了人类思维的不同侧面：逻辑、联想与因果。它们的交错演化，也构成了人工智能发展史中最富张力的部分。

\section*{符号主义：逻辑机器的黄金时代}

在人工智能的黎明时期（1950s–1980s），主流范式是符号主义，也称为“高层次AI”或“经典AI”。它试图用类似人类思维的方式去建模智能：使用规则、表达式、逻辑推理系统。

例如，专家系统（expert system）会将医生诊断的知识整理为“如果–那么”规则库：  
\texttt{if 症状 = 发烧 \& 咳嗽 then 疾病 = 流感}

早期的符号主义取得了惊人的成功，例如：
- 逻辑定理证明程序（如 Logic Theorist）
- 国际象棋程序（如 Deep Thought）
- 医疗诊断系统（如 MYCIN）

\textbf{优势}：
- 可解释性强，推理路径清晰；
- 易于集成专家知识；
- 可以精确表达语义和规则。

\textbf{局限性}：
- 对世界的表达非常“脆弱”：规则难以穷尽、难以处理模糊或例外；
- 缺乏学习能力：无法从数据中归纳新规则；
- 对于感知问题（如图像识别）无能为力。

随着现实世界复杂度的上升，纯符号系统开始暴露出“知识获取瓶颈”，也就是所谓的“知识工程灾难”。

\section*{连接主义：从神经元到深度学习}

神经网络的兴起，尤其是深度学习的突破，标志着连接主义的全面复兴。

连接主义的核心理念是：\textbf{智能并不是规则的堆砌，而是联想和模式的识别}。神经网络用大量训练数据调整权重，从而学习到复杂的输入–输出映射。

近年来的连接主义代表成果有：
- 图像识别：CNN 几乎主宰了整个视觉任务；
- 自然语言处理：Transformer 架构开启了语言建模新纪元；
- 再到 GPT 系列，甚至是多模态模型如 Gemini、Claude。

\textbf{优势}：
- 强大的感知能力；
- 可在大规模数据中发现复杂模式；
- 能适应非结构化、高维数据。

\textbf{局限性}：
- 黑箱结构，缺乏可解释性；
- 难以进行逻辑推理与常识判断；
- 泛化能力差，不擅长举一反三；
- 缺乏对“因果关系”的表达能力，只会“相关”而非“为什么”。

连接主义更像“右脑思维”：强感知、弱逻辑。而人类智能显然不止于此。

\section*{因果主义：预测之后，我们要理解}

在深度学习取得惊人成绩之后，部分学者开始反思：\textbf{预测能力 ≠ 理解能力}。AI 不仅要“看见模式”，还要“推理因果”。

因果主义代表人物裴珠叶（Judea Pearl）提出：真正的智能系统，应具备三层因果能力：

\begin{itemize}
\item \textbf{关联层}（Association）：观察现象之间的统计关系（如：下雨与打伞）；
\item \textbf{干预层}（Intervention）：如果我“做”了某件事，会发生什么？（如：我不打伞会不会湿？）；
\item \textbf{反事实层}（Counterfactual）：如果我当时打了伞，结果会不会不同？
\end{itemize}

深度学习大多停留在第一层，而真正理解世界的能力来自后两层。

因果推理的工具包括：
- 贝叶斯网络
- 因果图
- do-calculus 推理规则
- Rubin 的因果模型等

这些工具可以赋予模型某种“科学家的思维”：不是仅仅统计，而是构建假设、施加干预、分析结果。

\section*{神经-符号融合：一场重新组合的尝试}

随着符号主义和连接主义各自的局限暴露，研究者开始尝试一种“中间道路”：

\textbf{能否将神经网络的感知能力，与符号系统的逻辑推理能力结合？}

这就是近年来兴起的“神经–符号融合”（Neuro-Symbolic AI）研究方向。

典型的尝试包括：
- \textbf{Neuro-Symbolic Concept Learner (NS-CL)}：利用神经网络从图像中提取对象特征，用符号模块完成逻辑问答；
- \textbf{Logic Tensor Networks}：将逻辑规则嵌入神经网络中；
- \textbf{DeepProbLog}：在概率逻辑编程中加入神经网络模块；
- \textbf{AlphaCode}：使用语言模型生成代码，同时结合语义约束与程序验证。

这类模型具有以下特征：
- 感知与推理协同；
- 可解释性增强；
- 具备基本的归纳能力；
- 对知识结构与逻辑条件敏感。

神经–符号融合是一种对“通用智能”的重新构想：我们或许无法靠神经网络一口气吃下整个智能，也无法靠规则系统编码整个世界，但\textbf{两者合起来，或许足以搭建起更高阶的认知结构}。

\section*{未来：从三元对立到认知整合}

我们不妨这样看待：
- \textbf{符号主义}代表的是“结构化知识”；
- \textbf{连接主义}代表的是“模式联想”；
- \textbf{因果主义}代表的是“推理与干预”。

三者并非敌对，而是共同构成人类思维的不同侧面。真正的智能，不应是一场零和竞争，而是一种整合：\textbf{结构中的联想，联想中的因果，因果中的规则}。

未来的人工智能，也许不再强调“深度”，而更强调“层次”；不再执着于“神经网络参数的多少”，而是追问——这些参数里是否真的有“推理的种子”。

下一章，我们将关注人工智能在语言中的表现——它最擅长的领域，也是它最易暴露幻觉的舞台。Transformer 的崛起，是否真的让机器开始“理解”语言？

\nocite{*}
\printbibliography[heading=bibintoc, title=\ebibname]

\newpage
